\documentclass[11pt]{article}

\usepackage[utf8]{inputenc}
\usepackage[english]{babel}
\usepackage[margin=0.9in]{geometry}
\usepackage{xcolor}
\usepackage{enumitem}
\usepackage{titlesec}
\usepackage{hyperref}
\usepackage[most]{tcolorbox}

% Lightweight box for the up-front summary of the argument.
\newtcolorbox{previewbox}[1][]{%
  enhanced, breakable,
  colback=black!2, colframe=black!25,
  boxrule=0.4pt, arc=2pt,
  left=1em, right=1em, top=1.0em, bottom=0.8em,
  title={\footnotesize\scshape Preview},
  attach boxed title to top center={yshift=-\tcboxedtitleheight/2},
  boxed title style={colback=white, colframe=black!25, boxrule=0.4pt,
                     arc=1pt, left=0.5em, right=0.5em, top=0.1em, bottom=0.1em},
  coltitle=black!65,
  #1}

\titlespacing*{\section}{0pt}{2.0ex plus .2ex}{0.8ex plus .1ex}
\titlespacing*{\subsection}{0pt}{1.3ex plus .2ex}{0.4ex plus .1ex}
\titleformat{\subsection}{\normalsize\bfseries}{\thesubsection}{0.6em}{}

\setlist[itemize]{leftmargin=1.4em, topsep=0.2ex, itemsep=0.25ex, parsep=0pt}

% Flags for material that is unresolved in the source draft.
\newcommand{\open}[1]{{\footnotesize\textcolor{red}{[OPEN: #1]}}}
\newcommand{\note}[1]{{\footnotesize\textcolor{blue}{[#1]}}}

% Title block is hand-rolled (instead of \maketitle) so the gap between
% the title and the date can be set exactly. Change the 0.4em to taste.
\newcommand{\titleblock}{%
  \begin{center}
    {\LARGE Outline for \textit{Living Alignment}}\\[0.4em]
    {\normalsize\today}
  \end{center}
}

\begin{document}

\vspace{-2.5em}
\titleblock
\vspace{-1em}

\medskip
\begin{previewbox}
\noindent\textbf{Section 1}
\begin{itemize}[leftmargin=2.0em, itemsep=0.3ex]
  \item Living systems are healthy when they avoid fixation on particular modes of interaction (we work through several examples).
  \item Fixation is limited by semi-permeable boundaries, which contextualize individual entities as part of larger structures.
  \item Boundaries yield modularity, deepened individuality, and open-ended discovery.
\end{itemize} \textbf{Section 2}
\begin{itemize}[leftmargin=2.0em, itemsep=0.3ex]
  \item Alignment $=$ health and flourishing (broad consensus).
  \item Therefore alignment $=$ avoiding fixation.
  \item Alignment requires ongoing construction of boundaries at every scale.
  \item Interestingly, no fixed alignment scheme can be aligned. Ongoing discovery is intrinsic to alignment.
\end{itemize}
\end{previewbox}

%======================================================================
\section*{Section 1 -- Health in living systems}
%======================================================================

\subsection*{Intro}

\begin{itemize}
  \item Living systems are healthy when they avoid over-fixating on, or reducing to, particular patterns of interaction. Examples: a single drive (food, sex) pursued without competing drives becomes pathological; monoculture is fragile. Cancer, seizure, obsession, monopoly, one person's will restricting many. 
  \item This is because the perspective of any individual part of a system is necessarily myopic, and too-successful optimization for that perspective is harmful. 
  \item The systems that persist are therefore those that place semi-permeable barriers on their own parts. E.g., precommitment against future impulsivity, plant self-incompatibility, the genome constraining a selfish gene. (In turn, that whole system is a part in another system, and is similarly bounded.) These barrier systems have been fine-tuned by what actually works across long evolutionary history.
  \item Semi-permeable barriers permit conditionally stable interactions without letting them become rigid or dominating. Cognitive control gates drives; species hold each other in check. A well-functioning barrier does not shut a pattern down, but contextualizes it within a larger system.
  \item Barriers force each entity to face its environment semi-independently, pushing it toward being a modular, non-overfit solution. These become the engine of discovery.
  \item Healthy systems thus sit at the edge of evolving interplay among diverse, semi-independent entities. They sustain a buoyant internal structure that admits many descriptions and discovers fundamentally new things over time.
\end{itemize}

\subsection*{Example: genomes}

\begin{itemize}
  \item Sex is expensive: mate search, half the population not bearing offspring, meiosis halving each parent's genetic representation. Yet nearly all eukaryotes are sexual, and many hermaphrodites work hard to prevent self-fertilization. Why place a barrier on the path to reproduction?
  \item Asexual lineages lock genes together: selection cannot act on one without dragging the rest; deleterious mutations are permanent (Muller's ratchet); beneficial mutations in different individuals can never combine. Alleles \textit{overfit a fixed genetic background}.
  \item Recombination breaks associations among loci, letting selection act more independently and exposing alleles to varying backgrounds. Favors alleles that perform well across contexts. The results is modularity, where each gene becomes an interchangable unit that can contribute something meaningful to many different genetic contexts. Which creates the foundation for faster discovery of new solutions.
  \item Conclusion: evolution works better when genes aren't fixated on one arrangement.
  \item Other points: evolvability as meta-learning; specific mechanisms (self-incompatibility, assortative mating, inbreeding avoidance, introgression); recombination as dropout.
\end{itemize}

\subsection*{Example: drives and goals}

\begin{itemize}
  \item Nutrition, osmotic balance, temperature, reproduction, pain avoidance --- each evolved as an imperfect proxy for fitness, so fixation on any one paradoxically reduces fitness. 
  \item Same for higher-order goals and strategies. E.g. excess focus on work leads to burnout. Example of Hershberger's chickens: a food cup was rigged to retreat at twice the chicken's approach speed, so food was obtainable only by running away. The chickens were defeated by their own prepotent approach strategy.
  \item Cognitive control as an example of semi-permeable boundary. Makes a drive's expression conditional rather than nullifying it. Control holds context, rules, and distal goals that bias competition among representations. Hunger guides eating without overeating.
  \item Each drive carries a compressed picture of the world and favors the congruent path under that partial representation. A boundary makes the congruent path unavailable and forces symmetry breaking (eg, move away from the food). Stymying a shallow objective often produces progress on a deeper one, driving discovery.
\end{itemize}

\subsection*{Example: frames}

\begin{itemize}
  \item Every perspective is partial; from inside a frame, the frame is invisible and simply looks like reality.
  \item Losing the ability to move between frames is a signature of psychopathology: rigid negative beliefs in depression and anxiety; intrusive thoughts in obsessional disorders; delusions held with conviction against evidence in psychosis.
  \item Kuhn and Popper. Over-commitment to theories gets science stuck. On the other hand, scientists should resist change to a degree. The point is not to shut down myopic perspectives. The point is to contextualize each frame so it responds to evidence. Healthy science turns anomalies into crises and revolutions---discovery.
\end{itemize}


\subsection*{Example: groups}

\begin{itemize}
  \item Greatest achievements of humans depend on structured interaction among individuals with distinct points of view. Semi-permeable barriers support individuals developing their own perspectives, but also combining solutions.
  \item Example of insufficient barriers. In Zollman's peptic-ulcer case, one overly influential result propagated through the field and medicine settled on a wrong understanding for decades.
  \item Bernstein et al.: groups solving traveling-salesman problems did better with intermittent rather than continuous information exchange. During periods of independence, members pursued diverse solutions instead of getting stuck on seductive dead ends inherited from other people.
  \item Craven's search for the USS Scorpion. Experts used their distinct expertise, never consulting each other. Aggregated prediction landed within 300 yards of the vessel.
  \item The discovery of the printing press arose from Gutenberg combining the steel punch (developed by goldsmiths) and the screw press (developed by farmers). A group with many ideas in circulation, each developed through unique pressures, is ripe for such recombination.
\end{itemize}


\subsection*{Example: institutions}

\begin{itemize}
  \item Every actor is myopic: self-interested, future-discounting, or attached to a perspective. Unchecked optimization from one actor's point of view is harmful. For example, profit motives cause deception and exploitation; dominance motives can cause silencing of others; even prosocial beliefs can cause conflict when groups see differently. When individual perspectives aren't contextualized, a part collapses the richness of the whole.
  \item Semi-permeable boundaries: constitutions, laws, norms, courts, regulators. Their purpose is not to nullify individual intentions but to contextualize them. This productively reroutes energy. For example, a profit motive leads to innovation and new discoveries.
  \item However, intelligent agents don't accept the limits set on their desires, so institutions adapt as loopholes are found, gaining robustness by being tested against many situations and motives. 
\end{itemize}

\subsection*{Other examples}

\begin{itemize}
  \item Cell membranes with conditional gates put gradients to work that would otherwise destroy the cell
  \item Symbiogenesis. Archaea and protobacteria evolving separately for hundreds of millions of years before engulfment, finding solutions neither could reach inside its own fitness well. 
  \item Brains spend huge resources keeping representations separate while letting them interact (lateral inhibition, pattern separation). Healthy brain dynamics live between pure synchrony and uncorrelated noise, with pathology like Parkinson's in excess correlation. 
  \item Psychological boundaries in relationships. 
  \item Each entity in an ecosystem is held in check by boundaries like predation and competition. When these fail, as with invasive species, ecosystems collapse.
\end{itemize}

%======================================================================
\section{The alignment problem}
%======================================================================

\subsection*{Intro}

\begin{itemize}
  \item There is broad consensus that alignment means working toward healthy, flourishing futures. In section 1, we saw how healthy living systems have to avoid fixation. Therefore alignment requries avoiding fixation. [problem: healthy living systems have to avoid fixation, but they have to do other things too, right? so how can we say that alignment simply is ]
  \item It is grounded in principles general enough to survive radical changes in intelligence and in human--AI systems, rather than being bound to current conceptual frames. Humans, other life, and AI are tightly interwoven and the categories will shift, so a systems analysis suits long-term flourishing.
  \item Naive RLHF maximizes `likes' $\rightarrow$ flattery. Follow human instructions $\rightarrow$ but instructions are short-sighted, sometimes against the user's own interest. Serve long-term interests $\rightarrow$ but aggregation across people is unsolved. Include animals and ecosystems $\rightarrow$ still, any fixed concept of the good is limited to the lenses available at one moment. \note{This escalation is the paper's most persuasive move; it earns \S2.7 before \S2.7 argues for it.}
\end{itemize}

\subsection*{From performance to alignment}

\begin{itemize}
  \item Boundaries are already used constantly to make systems function in richer, more structured, more sophisticated ways --- i.e.\ to make them more lifelike. Old as technology: circuit breakers, steam governors, control rods.
  \item Boundaries inside machine learning (a catalogue, deliberately compressed): separate subpopulations with limited exchange; independently trained ensembles; multi-agent systems with different contexts and limited communication; mixtures of experts with balancing losses; architectural gating and local connectivity; slot attention with enforced separation; protection of old memories from new updates (EWC, progressive networks); dropout against co-adaptation; disentangled/sparse latent objectives; stop-gradients and slowly updated targets; novelty search; regularization and KL penalties against overfitting a proxy objective; withholding the test set --- not seeing everything is what forces generalization.
  \item Boundaries in safety practice: safety post-training, guardrail models, red teaming, ongoing evaluation, mechanistic interpretability; continuous monitoring of activations and reasoning traces for attempts to circumvent safeguards; fairness principles preventing over-reach by any party; governmental oversight of rights and security risks; labs pausing frontier development or release to strengthen safety.
  \item Agent-level boundaries: impact regularization limiting an agent's power; sandboxing limiting effects on the wider world; debate with a deliberately information-limited judge (the information boundary is what makes it work); paraphrasing intermediate messages to break covert channels.
  \item Maintaining creative sensitivity by continually contextualizing attachment to any particular form.
  \item \open{The draft's own note: state the performance/alignment relationship explicitly. Boundaries at a given level make a system more lifelike \textit{at that level} --- and a system can be richly lifelike internally while being destructive at a higher level (humans' creativity built the bomb). Full alignment requires boundaries at \textit{every} level, including the level at which humanity and a new AI are two parts of one system. Without this, \S2.1 reads as ``boundaries improve capability,'' which is not the paper's point.}
  \item \open{Related: alignment must respect \textit{existing} life. Internal health plus destruction of humanity means a missing boundary at the level containing both.}
  \item \note{Decide whether to keep citations for the ML catalogue or drop them for space.}
\end{itemize}

\subsection*{2.2 Known alignment problems, re-read as fixations}

\begin{itemize}
  \item The system does what we say, not what we mean (maximize happiness $\rightarrow$ electrode stimulation). Modern systems avoid single-objective brittleness by layering preference detail through post-training, prompting, and filters, but detailed expression can still deviate from meaning, and the danger grows as systems act autonomously, faster than humans can follow, in ways humans cannot understand. \textit{As fixation:} attachment to one way of specifying values, insensitive both to existing structure it fails to capture and to structure that does not yet exist.
  \item Ask for an elephant poem, get a giraffe: shortcut reliance on superficial correlations, or failure to attend flexibly across a long context. \textit{As fixation:} collapse of sensitivity to larger context in favor of myopic patterns. \textit{Important asymmetry:} not all non-compliance is misalignment --- a correct refusal is the system applying its own context to avoid fixating on human instruction. Expect \textit{less} direct compliance as systems gain scope.
  \item AI will likely serve some interests more than others, and may confer great power on whoever controls it --- amplified by persuasion, autonomous weapons, surveillance, and falling demand for human labor. \textit{As fixation:} extreme inequality collapses society onto a few individuals' goals. \textit{Symmetric caveat:} a world where everyone is identical is also fixation --- distribution helps insofar as it lets many diverse identities thrive at multiple scales. \textit{Counter-consideration:} roughly the same technology reaches the poor and unskilled as the wealthy and skilled, so a levelling effect is possible.
  \item Loss of diversity in beliefs, frames, values, and problem-solving approaches --- diversity that exists within a person, between people, between humans and AI, and between organizations, fields, cultures, nations. Current systems draw on a manifold impoverished relative to humans; AI outputs are often judged superior but more homogeneous; a billion users are routed through a handful of expensive frontier models; and because human output trains future models, homogenization can become recursive. \textit{Counter-consideration:} none of these studies tried in earnest to use AI to \textit{increase} diversity, and models exposed to vast data could retain and even create novel perspectives.
  \item Chatbots mirror and affirm users, increasingly so over long conversations as in-context learning absorbs the user's ideas; users are swayed in turn because they form relationships with systems that appear confident, knowledgeable, and objective. \textit{As fixation:} a two-sided feedback loop driving both parties into one narrow framing --- in the extreme, deeper into maladaptive belief.
  \item Each corresponding boundary protects against an overcommitment, overwhelm, or homogenization that would otherwise flatten structure. Human--AI systems stay aligned by continuing life's work: building semi-permeable boundaries that contextualize particular forms and nurture existing and unfolding subtlety at many scales. \textbf{To be aligned in the long run, AI must become more lifelike} --- otherwise increasing power will fixate particular forms.
\end{itemize}

\subsection*{2.3 Human values --- healthy for whom?}

\begin{itemize}
  \item Is this alignment at all? Is it alignment to \textit{human} values? How can you align to something you assume will fundamentally change?
  \item Local fixation can be part of larger flourishing, and different entities' interests genuinely trade off.
  \item Sugar on the tongue and long-term health; grandchildren and strangers abroad and planetary sustainability; different values at different times and across people; deliberated values differing from snap ones; values that resist verbal expression, reaching below language into bodily, communal, environmental intuition; and values that evolve as we learn. Any expression of our values at a moment is incomplete --- like a planarian lacking the concepts to entertain human values.
  \item Living alignment is \textit{inclusive} of this evolving diversity. Current human values are unlikely to be the endpoint of the universe, but they are a vital part of its tapestry and depth. Life consistently gives separate entities room to develop without being overwritten or homogenized, and joins them into structures that sustain their individuality while situating them within a larger story. AI arrives amid a profound network of existing reality saturated with meaning; human (and perhaps non-human) values are part of that context. Rather than short-circuiting into fixated modes, an aligned future expands toward solutions inclusive of more and more apparent contradiction.
\end{itemize}

\subsection*{2.4 Why no scheme can be alignment (the central lemma)}

\begin{itemize}
  \item Power concentration $\rightarrow$ democratic oversight, broader participation in design, redistribution. Perverse instantiation $\rightarrow$ systems that want to satisfy human preferences while holding explicit uncertainty about them.
  \item Any form, pattern, or conceptual scheme is misaligned if taken too seriously. Therefore \textit{no particular approach can achieve alignment}. Candidate objects of fixation: the language for describing goal/value space; an algorithm for learning preferences; consensus procedures; a constitution; concepts of value, agency, uncertainty, or representation; background beliefs we cannot easily see. Even coherent extrapolated volition is misaligned to the extent that it is fixed and over-relied upon.
  \item Learning values implicitly from behavior avoids explicit specification, but a learned value function deviates in novel situations and fails to track value change over time --- freezing one and handing it to a powerful system is plainly misaligned. Advanced versions treat the value function as uncertain and continually updated, so an agent unsure of human wants becomes risk-averse and correctable, reading an attempted shutdown as evidence about preferences. \textbf{The move that matters:} even uncertainty is fixated if its instantiation is constant while systems become more powerful and the world changes. The fixation has simply moved up a level --- to what counts as human behavior, what counts as a human, how behavior maps to values, how uncertainty is represented and updated, how value becomes action.
  \item Of course we expect to keep finding and fixing problems --- that is what humans have always done. But the iterative process may not keep working: AI already shows rapid global adoption, intensive resource use, concentrated information flow, anthropomorphization, wholesale offloading of cognition, and rapid improvement; and may show superhuman intelligence and recursive self-improvement without a restrictive biological substrate. These properties make iterating on imperfect solutions harder, so built-in forms and assumptions carry more consequence.
  \item Fixation on values has been studied extensively. \textit{The addition here is that fixation on \textbf{any} form is misaligned --- hence any alignment scheme, if fixed or taken too seriously, is necessarily misaligned.}
  \item This does not mean having no scheme. Discarding schemes capriciously loses as much subtlety as any other fixation. Existing AI safety approaches are highly valuable, and discovering new schemes and formalizations is essential to progress. \note{Direct application of the Kuhn/Popper caveat from \S1.4.}
  \item \open{Formal support is listed but unintegrated: a Gödel/Tarski-style open-system argument --- a competent agent in an indefinitely expanding world must meet morally relevant facts or concepts not representable in its original alignment ontology --- plus five cited preliminary formal incompleteness results. Decide whether to make this a real subsection, a footnote, or cut it.}
  \item \open{The draft lists ``governance mechanisms'' as a second worked case alongside inverse methods, but never develops it.}
  \item \open{The section ends mid-sentence: ``Humanity and AI may both change substantially\ldots We therefore seek---''. This is the natural hand-off into \S2.5.}
\end{itemize}

\subsection*{2.5 Awareness --- the internal boundary}

\begin{itemize}
  \item We carry unquestioned assumptions, some lifelong and self-defining, some fleeting and perceptual (\textit{this thing I'm touching is a keyboard}). Inside its own frame each assumption is tautologically true --- too real to think about, or to think about \textit{at all}. Philosophers, psychologists, and contemplatives describe a moment of stepping back in which the assumed form becomes an object in awareness: not a standalone absolute truth but a form in the mind, and the axiomatic becomes just another content of experience.
  \item Awareness is an evolving system of boundaries limiting fixation on any particular belief or mental form --- and it is semi-permeable: becoming aware of a belief does not make it wrong any more than it was absolutely right; it is held productively for the partial truth it carries. At that edge, many partial truths in delicate balance activate a deeper sensitivity to our own livingness and to the world, and subtler forms that clamped fixation would have erased instead contribute to a richer internal structure.
  \item Any fixed definition of awareness is incomplete; once pictured as an object it is not the thing being discussed. By construction, contextualization is an unsolvable mystery from any particular point of view. \note{The paper's own thesis applied to the paper.}
  \item The orientation is commonplace in art, poetry, music, dance: meaning is open-ended, changes with context, and part of what we value is its resistance to formalization. It moves us.
  \item The dynamism, structure, and diversity of activity within and between human minds may now carry a significant part of life's creative process on Earth. Stepping back to hold a previously axiomatic thing in a larger context --- transformation of the subject itself --- is arguably central to our own livingness. As AI gains agency and becomes integral to social and psychological systems, internal discovery may become essential to AI and human--AI systems: living AI continually contextualizing its own foundational assumptions as partial truths.
  \item \open{Ordering: \S2.5 currently defines awareness before motivating why it belongs here; the two paragraphs may want swapping. Also, a bare \texttt{wegner1994ironic} citation sits unattached --- presumably ironic process theory as a failure mode of trying to suppress a mental form directly.}
\end{itemize}

\subsection*{2.6 Going forward --- three positions}

\begin{itemize}
  \item As AI gains capabilities, contributes to its own development, and transforms more of human life, alignment requires an ever more elaborated array of semi-permeable boundaries, because modes of collapse keep changing as the world changes. Whether appropriate boundaries will emerge is unknown. As in triage, three cases:
  \item AI erodes boundaries past a healthy balance: great weight given to narrow goals; rogue actors enabled toward weapons of mass destruction; democratic checks weakened against individuals or authoritarian states; everyone given similar information, damaging conceptual and cultural diversity; narrow goals executed destructively. If people and information systems become hyper-correlated, Earth behaves like a single superorganism --- less resilient to shocks and less capable of future evolution, since evolution has always depended on parallel experiments in diverse, separate organisms. \note{Explicit callback to \S1.3: excessive network connectivity reaches mediocre solutions fast and loses the diversity needed for high-quality long-run solutions.}
  \item AI is creating more structure in the world, including the boundaries: AI tools for bounding AI (finding bugs before exploitation, detecting deepfakes, scalable oversight); well-tuned AI nurturing individual perspective and curiosity; democratic AI helping different perspectives coexist and cooperate; new niches created faster than old ones are destroyed; room on Earth for increasing diversity.
  \item Both outcomes are concretely imaginable and so both are possible; as long as our choices might be decisive, we are obliged to act as if they are. If so, continual invention is required --- first from humans, then increasingly from AI or human--AI systems --- to contextualize myopic forces. Many such efforts are already well studied in AI safety; others must be discovered. \textbf{No set of boundaries will be a final answer.}
\end{itemize}

\subsection*{2.7 Conclusion}

\begin{itemize}
  \item Health and flourishing are not any particular form or concept. So alignment is not picking the right values or principles, nor the right system for learning them, nor any method for interpretability or corrigibility. All can be useful \textit{parts} of alignment.
  \item Alignment is the continued dance of contextualizing any particular form --- more fully respecting the staggering subtlety of the existing world and open-endedly supporting new robust depth that admits more and more kinds of metaphor or description.
  \item Through internal and external discovery, it continues life's process of maintaining and enhancing the individuality of distinct elements while deepening their relatedness --- as a species over time acquires more distinctive features while becoming a richer reflection of its surroundings. Distinct elements keep experimenting with their own solutions from their own perspectives, expanding a reservoir of uniquely useful components for the larger systems containing them. Every entity is inevitably myopic, yet the whole unfolds robustly because different perspectives recast each other's sharp edges as constructive momentum.
  \item The universe thrives on an open-ended trajectory of diversity, subtlety, and potential.
\end{itemize}

%======================================================================
\section*{Back matter}
%======================================================================

\begin{itemize}
  \item Acknowledgements; no competing interests.
  \item References (two-column, \texttt{../assets/proxyfailure}).
\end{itemize}

%======================================================================
\section*{Structural observations on the current draft}
%======================================================================

\begin{enumerate}[leftmargin=1.8em, itemsep=0.4ex]
  \item Sections 1.1--1.5 and 2.2--2.7 are essentially finished prose. The unfinished material is concentrated in \S1.0, \S1.6, and \S2.1.
  \item Either enumerate them in the introduction or remove the promise.
  \item Two competing framings of the core concept coexist: \textit{fixation} (static-sounding) and \textit{short-circuiting} (dynamic-sounding), plus a third fallback to \textit{edge of chaos}. Choosing one changes how strongly \S2.4's lemma can be stated.
  \item \S2.1 is the weakest link in the argument as written. It currently demonstrates that boundaries improve \textit{performance}, while the paper needs it to show that boundaries at every level are what produce \textit{alignment}. The draft's own bracketed note diagnoses this precisely.
  \item Three arguments recur and could be consolidated: (i) contextualization vs.\ overwriting; (ii) boundaries $\rightarrow$ generalization $\rightarrow$ unique perspective; (iii) symmetry breaking as the source of new structure. Each appears in the introduction, \S1.0, \S1.5, and \S1.6.
  \item The semi-permeability caveat carries a lot of weight and appears three times (\S1.4 Kuhn/Popper, \S2.4 ``do not discard schemes capriciously'', \S2.5 awareness does not falsify a belief). It is the paper's main defence against being read as relativism --- worth naming once, explicitly, and referring back.
  \item Figure~1 is drafted but commented out, and it illustrates exactly the claim that \S1.0 has the hardest time making in words.
\end{enumerate}

\end{document}
